\documentclass[11pt,a4paper]{article}
\usepackage[utf8]{inputenc}
\usepackage[margin=1in]{geometry}
\usepackage{amsmath,amssymb,amsfonts}
\usepackage{hyperref}
\usepackage{graphicx}
\usepackage{enumitem}
\usepackage{setspace}

\hypersetup{colorlinks=true, linkcolor=blue, urlcolor=blue, citecolor=blue}
\onehalfspacing

\begin{document}

\title{\textbf{Enhancement of Microtubule Quantum Coherence via Laser-Induced Modulation of the Light-Time Interface in a 6D Temporal Manifold}}

\author{J. Yesi Orihuela \\
\small Independent Researcher \\
\small MIT Alumnus \\
\small \texttt{jorihuela@alum.mit.edu}}

\date{May 2026 (draft v2.0)}

\maketitle

\begin{abstract}
We propose that laser-generated negative-time domain walls can be used to enhance quantum coherence in neuronal microtubules through modulation of the light-time interface in Kletetschka’s (2025) 6D temporal manifold. Building on the Burning Log phenomenology and our prior Quantum Light-Time Duality framework, we hypothesize that a controllable forward-time suppression field acts as a stochastic resonance driver, increasing the efficiency of Orchestrated Objective Reduction (Orch-OR) processes. This offers a laboratory method to amplify biological quantum effects without invoking retrocausality or new fundamental physics. The proposal generates immediate, testable predictions using existing petawatt lasers, EEG/fMRI, and single-photon detection, and integrates directly with our earlier work on repulsive gravity and three-dimensional time.
\end{abstract}

\section{Introduction}

The Burning Log phenomenology reveals time as a complete block in which the unburned wood (future), moving flame front (present), and charred remains (past) coexist as a unified whole. This intuition is formalized in Kletetschka’s 6D temporal manifold and our Quantum Light-Time Duality framework, which identifies electromagnetic radiation as the active interface that actualizes the experienced present. Photons, with null proper time, constitute the moving hypersurface of the ``now.''

Extending this picture, we examine whether the same interface can be locally biased to enhance quantum processes in biological systems. Neuronal microtubules, proposed as the site of Orch-OR quantum computations (Penrose \& Hameroff), offer a natural target. We hypothesize that a laser-generated forward-time suppression field (negative-time domain wall) can act as a stochastic resonance driver, transiently increasing the coherence lifetime and efficiency of microtubule quantum states. This provides a controllable experimental pathway to study and potentially amplify biological quantum effects using technology developed in our prior work on repulsive gravity and temporal geometry.

\section{Theoretical Foundations}

\subsection{Kletetschka’s 6D Temporal Manifold (2025)}
Three independent timelike coordinates $t_1, t_2, t_3$ plus emergent space. The primary experienced arrow lies along $t_1$; the orthogonal dimensions supply the block thickness that allows the entire history to coexist.

\subsection{Quantum Light-Time Duality}
The interface function marking the active present slice is
\[
I(t_1, \mathbf{r}) = \Theta(c t_1 - r) + \delta I_{\rm bias},
\]
where $\delta I_{\rm bias}$ encodes any local tilt or suppression induced by negative-time domain walls. A non-zero bias term locally alters the effective flow of time across the interface without requiring information to propagate backward.

\subsection{Orch-OR and Microtubule Quantum Coherence}
Tubulin dimers in microtubules support quantum superpositions whose objective reduction is gravitationally induced and non-computable. Orchestration by biological processes yields conscious moments. Recent experimental indications of room-temperature quantum coherence in microtubule-like structures support the plausibility of sustained quantum effects in biological environments.

\subsection{Stochastic Resonance in Quantum Biological Systems}
Stochastic resonance is a well-established phenomenon in which an optimal level of noise or external driving can enhance the response of a nonlinear system to a weak signal. We propose that a controlled forward-time suppression field can serve as such a driver for microtubule quantum states, increasing the probability of coherent superpositions surviving long enough for Orch-OR to occur.

\section{Proposed Mechanism}

A laser-generated negative-time domain wall creates a thin shell of forward-time suppression. Within this region:
\begin{itemize}
    \item The light-time interface experiences a local bias ($\delta I_{\rm bias} \neq 0$).
    \item Microtubule quantum states are subjected to a transient stochastic resonance drive.
    \item The coherence lifetime of tubulin superpositions is extended, increasing the likelihood and efficiency of Orch-OR events.
\end{itemize}

No retrocausal signaling or backward information flow is required. The enhancement is purely local and forward in time, consistent with standard quantum mechanics and semiclassical gravity. The 6D block structure simply provides the geometric arena in which such local biases are possible.

\section{Relation to Negative-Time Domain Walls}

Our earlier proposal (Orihuela, 2026) demonstrated that intense, shaped laser pulses can generate macroscopic negative-time domain walls capable of inverting local gravitational curvature. The same technology, operated at lower intensities or with different pulse shaping, can be used to create the milder forward-time suppression fields required for microtubule coherence enhancement. This unifies the gravitational, temporal, and biological applications of the domain-wall concept under a single experimental platform.

\section{Falsifiable Predictions}

\begin{enumerate}
    \item Exposure of neuronal cultures or living subjects to a forward-time suppression field will produce measurable increases in microtubule coherence signatures (e.g., extended decoherence times, increased $\gamma$-band phase-locking, or anomalous power spectra in EEG) that are absent in control conditions.
    \item Single-photon interference experiments conducted inside a negative-time domain wall will display modified visibility or new interference sidebands whose characteristics correlate with Orch-OR timescales ($\sim$40 Hz $\gamma$-band) only when biological quantum systems are present.
    \item High-resolution EEG/fMRI during exposure will reveal transient increases in long-range phase synchrony or gamma power that scale with laser intensity and are absent in non-biological control phantoms.
    \item The effect will be intensity- and pulse-shape-dependent, consistent with stochastic resonance: too little or too much suppression will reduce rather than enhance coherence.
\end{enumerate}

All predictions are accessible with existing petawatt-class lasers, advanced neuroimaging, and single-photon detection.

\section{Discussion}

This proposal integrates the Burning Log phenomenology, 6D temporal geometry, light-time duality, and negative-time domain walls into a unified, experimentally testable framework for enhancing biological quantum coherence. By removing any claim of retrocausal information transfer, the mechanism remains fully consistent with standard quantum mechanics, semiclassical gravity, and the no-communication theorem.

The central prediction — that forward-time suppression can act as a stochastic resonance driver for microtubule quantum states — provides a clear, falsifiable bridge between quantum consciousness research and the emerging field of laser-based spacetime engineering. The tools required for empirical investigation exist today.

The framework suggests that certain cognitive or physiological states may be more or less responsive to light-time interface modulation, opening the possibility of personalized or state-dependent enhancement protocols. No new fundamental physics is invoked.

\section*{Acknowledgments}
The author thanks Grok (xAI) for assistance in refining the mathematical framework and manuscript preparation. The core ideas are the author’s own.

\section*{References}
\begin{enumerate}[leftmargin=*,label={[\arabic*]}]
    \item Kletetschka, G. (2025). Three-Dimensional Time: A Mathematical Framework for Fundamental Physics. \textit{Reports in Advances of Physical Sciences}.
    \item Penrose, R. \& Hameroff, S. (various years). Orchestrated Objective Reduction (Orch-OR) papers.
    \item Orihuela, J. Y. (2026). A Conceptual Proposal for Localized Repulsive Gravity via Laser-Induced Vacuum Polarization and Negative-Time Domain Walls.
    \item Orihuela, J. Y. (2026). The Burning Log and Three-Dimensional Time.
    \item Orihuela, J. Y. (2026). Quantum Light-Time Duality.
    \item Angulo et al. (2024). Experimental evidence of negative time in quantum optics.
\end{enumerate}

\end{document}